% Локальный макрос: кириллическая переменная в math-режиме.
% Math-режим в XeLaTeX берёт глифы из math-шрифта, в котором кириллицы нет;
% \rv{...} оборачивает символ в \text{}, что использует основной текстовый
% шрифт (Times New Roman) и корректно отображает русские буквы.
\providecommand{\rv}[1]{\text{#1}}

\section{Экономическое обоснование проекта по разработке программного обеспечения}

Цель настоящего раздела~--~подтвердить экономическую целесообразность разработки и применения системы компьютеризации ввода и мониторинга информации об абитуриентах при поступлении в БНТУ.
Раздел охватывает четыре содержательных вопроса: описание функций, назначения и потенциальных пользователей программного продукта, расчёт затрат на разработку, оценку экономического эффекта от внедрения и расчёт показателей эффективности инвестиций.

\subsection{Описание функций, назначения и потенциальных пользователей программного обеспечения}

Разработанное программное обеспечение представляет собой веб-приложение для автоматизации работы приёмной комиссии Белорусского национального технического университета.
Программный продукт предназначен для замены применяемых в настоящее время разрозненных средств учёта (электронных таблиц с макросами, локальных баз данных, ручного пересчёта рейтингов) единой централизованной информационной системой.

Основными функциями программного обеспечения являются:
\begin{itemize}
\item ведение организационной структуры университета~--~факультетов, кафедр, специальностей~--~через интерфейс администратора;
\item ведение конкурсных списков и категорий приёма с настраиваемыми планами, квотами и приоритетами защиты;
\item ведение справочника показателей оценивания абитуриентов и групп критериев с настраиваемыми правилами ранжирования;
\item ввод данных абитуриентов и их заявок в категории приёма;
\item автоматический пересчёт результатов отбора при любом изменении данных абитуриентов;
\item валидация введённых данных силами второго оператора по принципу двойного контроля;
\item ведение журнала действий пользователей (аудит) и журнала авторизации;
\item двухступенчатое удаление абитуриентов (soft-delete) с возможностью отката;
\item выгрузка результатов отбора в формате Excel для последующей обработки.
\end{itemize}

Потенциальными пользователями системы являются:
\begin{itemize}
\item сотрудники приёмной комиссии БНТУ, ответственные за ввод и обработку данных абитуриентов;
\item сотрудники приёмной комиссии БНТУ, выполняющие функции аудиторов и подтверждающие корректность введённых данных;
\item ответственные сотрудники кафедр и факультетов, ведущие справочники структуры и конкурсных списков;
\item администраторы информационной системы БНТУ, выполняющие конфигурирование, управление учётными записями и сопровождение;
\item опосредованно~--~абитуриенты, получающие достоверные и оперативно публикуемые результаты конкурсного отбора.
\end{itemize}

Потребность в программном продукте обусловлена ежегодным масштабом приёмной кампании БНТУ: количество подаваемых заявлений составляет порядка пяти тысяч и более, при этом нормативная база приёма периодически меняется, а ручной пересчёт рейтингов с использованием электронных таблиц трудоёмок и подвержен ошибкам.
Внедрение системы позволит сократить трудоёмкость рутинных операций, повысить достоверность результатов отбора за счёт автоматического пересчёта и валидации, а также обеспечить полноту прослеживаемости действий пользователей и истории изменений сущностей.

Косвенно потенциальными пользователями являются другие учреждения высшего образования Республики Беларусь, использующие сходную модель конкурентного зачисления по приоритетам и нуждающиеся в инструменте, который не привязан к нормативной базе других государств и допускает изменение правил приёма через интерфейс администратора.

\subsection{Расчёт затрат на разработку программного обеспечения}

Затраты на разработку программного обеспечения определяются по статьям: основная заработная плата команды разработчиков; дополнительная заработная плата; отчисления на социальные нужды; прочие затраты.

\subsubsection{Затраты на основную заработную плату команды разработчиков}

Затраты на основную заработную плату определяются исходя из состава и численности команды разработчиков, размеров месячной заработной платы каждого участника и общей трудоёмкости разработки программного обеспечения и рассчитываются по формуле

\begin{equation}
\label{eq:zo}
\rv{З}_\rv{о} = \rv{К}_\rv{пр} \cdot \sum_{i=1}^{n} \rv{З}_{\rv{ч},i} \cdot t_i,
\end{equation}

\begin{explanationx}
\item где $\rv{К}_\rv{пр}$~--~коэффициент премии, принимаемый в диапазоне 1,5--2,0;
\item $\mathrm{n}$~--~количество исполнителей, занятых разработкой;
\item $\rv{З}_{\rv{ч},\mathrm{i}}$~--~часовая заработная плата $\mathrm{i}$-го исполнителя, руб.;
\item $\mathrm{t}_\mathrm{i}$~--~трудоёмкость работ, выполняемых $\mathrm{i}$-м исполнителем, ч.
\end{explanationx}

Часовая заработная плата $i$-го исполнителя определяется делением месячного оклада на среднемесячный фонд рабочего времени, принимаемый равным 168~ч.
Размеры месячной заработной платы приняты на уровне рыночных значений по данным портала rabota.by для города Минска по состоянию на 2026~г.
Коэффициент премии принят равным $\rv{К}_\rv{пр} = 1{,}5$.

Состав команды разработчиков, трудоёмкость работ и расчёт затрат на основную заработную плату приведены в таблице~\ref{tab:zo}.

\begin{longtblr}[
  caption = {Расчёт затрат на основную заработную плату команды разработчиков},
  label = {tab:zo},
]{
  colspec = {|X[0.4,c,m]|X[1.9,l,m]|X[2.4,l,m]|X[1,c,m]|X[1,c,m]|X[2.1,c,m]|X[1.2,c,m]|},
  cells = {font=\small, cmd={\sloppy}},
  leftsep = 2pt,
  rightsep = 2pt,
  width = \textwidth,
  hlines, vlines,
  rowhead = 1,
  row{1} = {halign=c},
}
№ & Участник команды & Вид выполняемой работы & Оклад, руб. & Ставка, руб./ч & Трудоёмкость, ч & З/п по тарифу, руб. \\
1 & Бизнес-аналитик & Анализ требований, проектирование доменной модели & 2400 & 14,29 & 60 & 857,40 \\
2 & Full-stack .NET-разработчик & Реализация серверной (ASP.NET Core, PostgreSQL) и клиентской (React) частей & 3500 & 20,83 & 480 & 9998,40 \\
3 & Тестировщик & Функциональное и регрессионное тестирование & 2400 & 14,29 & 120 & 1714,80 \\
\SetCell[c=6]{l} Итого по тарифу & & & & & & 12~570,60 \\
\SetCell[c=6]{l} Премия (50\%) & & & & & & 6285,30 \\
\SetCell[c=6]{l} Итого $\rv{З}_\rv{о}$ & & & & & & 18~855,90 \\
\end{longtblr}

\subsubsection{Затраты на дополнительную заработную плату команды разработчиков}

Дополнительная заработная плата включает выплаты, предусмотренные законодательством о труде (оплата трудовых отпусков, льготных часов, времени выполнения государственных обязанностей и других выплат, не связанных с основной деятельностью исполнителей), и определяется по формуле

\begin{equation}
\label{eq:zd}
\rv{З}_\rv{д} = \frac{\rv{З}_\rv{о} \cdot \rv{Н}_\rv{д}}{100},
\end{equation}

\begin{explanationx}
\item где $\rv{З}_\rv{о}$~--~затраты на основную заработную плату, руб.;
\item $\rv{Н}_\rv{д}$~--~норматив дополнительной заработной платы, \%.
\end{explanationx}

Норматив принимаем $\rv{Н}_\rv{д} = 15$~\%.
Тогда

\[
\rv{З}_\rv{д} = \frac{18\,855{,}90 \cdot 15}{100} \approx 2828{,}39 \text{ (руб.).}
\]

\subsubsection{Отчисления на социальные нужды}

Отчисления в Фонд социальной защиты населения и на обязательное страхование определяются в соответствии с действующими законодательными актами Республики Беларусь по формуле

\begin{equation}
\label{eq:rsoc}
\rv{Р}_\rv{соц} = \frac{(\rv{З}_\rv{о} + \rv{З}_\rv{д}) \cdot \rv{Н}_\rv{соц}}{100},
\end{equation}

\begin{explanationx}
\item где $\rv{Н}_\rv{соц}$~--~суммарный норматив отчислений на социальные нужды, \%.
\end{explanationx}

В соответствии с действующим законодательством Республики Беларусь принимаем $\rv{Н}_\rv{соц} = 34$~\%.
Тогда

\[
\rv{Р}_\rv{соц} = \frac{(18\,855{,}90 + 2828{,}39) \cdot 34}{100} \approx 7372{,}66 \text{ (руб.).}
\]

\subsubsection{Прочие затраты}

Прочие затраты (амортизационные отчисления, расходы на электроэнергию, аренду помещений и оборудования, расходы на управление) включаются в себестоимость разработки в процентах от затрат на основную заработную плату и определяются по формуле

\begin{equation}
\label{eq:zpz}
\rv{З}_\rv{пз} = \frac{\rv{З}_\rv{о} \cdot \rv{Н}_\rv{пз}}{100},
\end{equation}

\begin{explanationx}
\item где $\rv{Н}_\rv{пз}$~--~норматив прочих затрат, \%.
\end{explanationx}

Норматив принимаем $\rv{Н}_\rv{пз} = 100$~\%.
Тогда

\[
\rv{З}_\rv{пз} = \frac{18\,855{,}90 \cdot 100}{100} = 18\,855{,}90 \text{ (руб.).}
\]

Полная сумма затрат на разработку программного обеспечения находится суммированием всех рассчитанных статей затрат и приведена в таблице~\ref{tab:zr}.

\begin{longtblr}[
  caption = {Затраты на разработку программного обеспечения},
  label = {tab:zr},
]{
  colspec = {|X[l,m]|X[1,c,m]|},
  width = \textwidth,
  hlines, vlines,
  rowhead = 1,
  row{1} = {halign=c},
}
Статья затрат & Сумма, руб. \\
Основная заработная плата команды разработчиков & 18~855,90 \\
Дополнительная заработная плата команды разработчиков & 2828,39 \\
Отчисления на социальные нужды & 7372,66 \\
Прочие затраты & 18~855,90 \\
Общая сумма затрат на разработку $\rv{З}_\rv{р}$ & 47~912,85 \\
\end{longtblr}

Таким образом, общая сумма затрат на разработку программного обеспечения составляет 47~912,85~руб.

\subsection{Оценка экономического эффекта от внедрения программного обеспечения}

Разрабатываемое программное обеспечение предназначено для удовлетворения собственных потребностей БНТУ как государственного учреждения высшего образования и не предполагает свободной реализации на рынке.
В связи с этим экономический эффект от его применения состоит в экономии трудозатрат при выполнении операций приёмной кампании.

Источники экономии трудозатрат и их оценочные значения приведены в таблице~\ref{tab:savings}.
Часовая ставка оператора приёмной комиссии БНТУ принята равной 12~руб./ч.

\begin{longtblr}[
  caption = {Оценка экономии трудозатрат при внедрении системы},
  label = {tab:savings},
]{
  colspec = {|X[l,m]|X[1,c,m]|X[1,c,m]|},
  width = \textwidth,
  hlines, vlines,
  rowhead = 1,
  row{1} = {halign=c},
}
Операция & Экономия времени, ч/год & Экономия, руб./год \\
Ввод данных абитуриентов в структурированные формы (5000~заявлений по 10~мин) & 833 & 9996 \\
Кросс-проверка и валидация введённых данных (5000~заявлений по 3~мин) & 250 & 3000 \\
Автоматический пересчёт результатов отбора после изменений (50~пересчётов по 4~ч) & 200 & 2400 \\
Автоматическая выгрузка отчётов в Excel (30~выгрузок по 2~ч) & 60 & 720 \\
Снижение трудозатрат на исправление ошибок благодаря средствам аудита, валидации и soft-delete & 100 & 1200 \\
Итого экономия $\rv{Э}_\rv{з}$ & 1443 & 17~316 \\
\end{longtblr}

С учётом округления принимаем $\rv{Э}_\rv{з} \approx 17\,300$~руб./год.

Прирост текущих затрат, связанных с эксплуатацией системы ($\Delta\rv{З}_\rv{тек}$), включает дополнительные расходы на электроэнергию и амортизацию серверной нагрузки (около 300~руб./год), а также сопровождение и мелкие доработки силами ИТ-отдела БНТУ (около 700~руб./год).
Итого $\Delta\rv{З}_\rv{тек} \approx 1000$~руб./год.

Прирост чистой прибыли от внедрения программного обеспечения для собственных нужд определяется по формуле

\begin{equation}
\label{eq:dpch}
\Delta\rv{П}_\rv{ч} = (\rv{Э}_\rv{з} - \Delta\rv{З}_\rv{тек}) \cdot (1 - \rv{Н}_\rv{п}),
\end{equation}

\begin{explanationx}
\item где $\rv{Э}_\rv{з}$~--~экономия текущих затрат, полученная в результате применения программного обеспечения, руб.;
\item $\Delta\rv{З}_\rv{тек}$~--~прирост текущих затрат, связанных с эксплуатацией программного обеспечения, руб.;
\item $\rv{Н}_\rv{п}$~--~ставка налога на прибыль в соответствии с действующим законодательством Республики Беларусь, в долях единицы.
\end{explanationx}

Белорусский национальный технический университет как государственное учреждение высшего образования освобождён от уплаты налога на прибыль на доходы от основной деятельности (Налоговый кодекс Республики Беларусь, ст.~181) [\ref{src:laws:nalogkodex}], поэтому $\rv{Н}_\rv{п} = 0$.
Тогда

\[
\Delta\rv{П}_\rv{ч} = (17\,300 - 1000) \cdot (1 - 0) = 16\,300 \text{ (руб./год).}
\]

Полученный годовой экономический эффект составляет 16~300~руб./год.

\subsection{Расчёт показателей эффективности инвестиций в разработку программного обеспечения}

Затраты на разработку программного обеспечения $\rv{З}_\rv{р} = 47\,912{,}85$~руб. превышают годовой экономический эффект $\Delta\rv{П}_\rv{ч} = 16\,300$~руб./год, поэтому простая окупаемость составляет более одного года и экономическая целесообразность инвестиций оценивается на основе расчёта дисконтированных показателей: чистого дисконтированного дохода (ЧДД), срока окупаемости ($\rv{Т}_\rv{ок}$) и рентабельности инвестиций ($\rv{Р}_\rv{и}$).

Расчётный период принят равным четырём годам и включает один год разработки и три года эксплуатации программного обеспечения.

Поскольку экономический эффект и инвестиции относятся к разным годам расчётного периода, для обеспечения сопоставимости они приводятся к началу расчётного периода путём умножения на коэффициент дисконтирования соответствующего года~$t$, определяемый по формуле

\begin{equation}
\label{eq:alpha}
\alpha_t = \frac{1}{(1 + \rv{Е}_\rv{н})^t},
\end{equation}

\begin{explanationx}
\item где $\rv{Е}_\rv{н}$~--~норма дисконта, в долях единицы;
\item $\mathrm{t}$~--~порядковый номер года расчётного периода.
\end{explanationx}

Норма дисконта принимается на уровне ставки рефинансирования Национального банка Республики Беларусь, действующей на момент проведения расчётов, с добавлением двух процентных пунктов.
Ставка рефинансирования Национального банка Республики Беларусь, действующая с 19~марта 2025~г., составляет 9,5~\% [\ref{src:nbrb-rate}].
С учётом надбавки в два процентных пункта принимаем $\rv{Е}_\rv{н} = 11{,}5$~\%.

Чистый дисконтированный доход за расчётный период определяется по формуле

\begin{equation}
\label{eq:npv}
\rv{ЧДД} = \sum_{t=1}^{n} (\rv{Р}_t \cdot \alpha_t - \rv{З}_t \cdot \alpha_t),
\end{equation}

\begin{explanationx}
\item где $\mathrm{n}$~--~расчётный период, лет;
\item $\rv{Р}_\mathrm{t}$~--~результат (экономический эффект), полученный в году~$\mathrm{t}$, руб.;
\item $\rv{З}_\mathrm{t}$~--~затраты (инвестиции), осуществлённые в году~$\mathrm{t}$, руб.
\end{explanationx}

Рентабельность инвестиций рассчитывается как отношение суммы дисконтированных результатов к сумме дисконтированных инвестиций по формуле

\begin{equation}
\label{eq:pi}
\rv{Р}_\rv{и} = \frac{\sum_{t=1}^{n} \rv{Р}_t \cdot \alpha_t}{\sum_{t=1}^{n} \rv{З}_t \cdot \alpha_t}.
\end{equation}

Расчёт показателей эффективности инвестиционного проекта приведён в таблице~\ref{tab:npv}.

\begin{longtblr}[
  caption = {Расчёт эффективности инвестиционного проекта по разработке программного обеспечения},
  label = {tab:npv},
]{
  colspec = {|X[2,l,m]|X[1,c,m]|X[1,c,m]|X[1,c,m]|X[1,c,m]|},
  width = \textwidth,
  hlines, vlines,
  rowhead = 1,
  row{1} = {halign=c},
}
Показатель & 1~год & 2~год & 3~год & 4~год \\
Коэффициент дисконтирования $\alpha_t$ & 0,897 & 0,804 & 0,721 & 0,647 \\
Экономический эффект $\rv{Р}_t$, руб. & 16~300 & 16~300 & 16~300 & 16~300 \\
Дисконтированный результат $\rv{Р}_t \cdot \alpha_t$, руб. & 14~621 & 13~105 & 11~752 & 10~546 \\
Инвестиции $\rv{З}_t$, руб. & 47~913 & 0 & 0 & 0 \\
Дисконтированные инвестиции $\rv{З}_t \cdot \alpha_t$, руб. & 42~978 & 0 & 0 & 0 \\
ЧДД по годам, руб. & --28~357 & 13~105 & 11~752 & 10~546 \\
ЧДД нарастающим итогом, руб. & --28~357 & --15~252 & --3500 & 7046 \\
\end{longtblr}

По результатам расчёта получены следующие показатели эффективности инвестиционного проекта:
\begin{itemize}
\item чистый дисконтированный доход на конец расчётного периода составляет $\rv{ЧДД} = 7046$~руб., что больше нуля и подтверждает экономическую эффективность инвестиций;
\item срок окупаемости проекта находится между третьим и четвёртым годом расчётного периода и составляет ориентировочно $\rv{Т}_\rv{ок} \approx 3{,}33$~года;
\item рентабельность инвестиций составляет $\rv{Р}_\rv{и} = 50\,024 / 42\,978 \approx 1{,}164$, что больше единицы и также подтверждает экономическую целесообразность вложения средств в разработку.
\end{itemize}

Все три показателя согласованно подтверждают экономическую эффективность проекта.
При этом выполненный расчёт является консервативным: норма дисконта принята с надбавкой в два процентных пункта к ставке рефинансирования Национального банка Республики Беларусь, а в составе экономии текущих затрат учтены только прямые сокращения трудозатрат сотрудников приёмной комиссии; косвенные эффекты внедрения~--~снижение количества ошибок ручного ввода, сокращение времени получения промежуточных рейтингов абитуриентами, повышение прозрачности приёмной кампании~--~в денежной форме не оценивались.

Срок окупаемости 3,33~года формально превышает половину расчётного периода, однако согласуется со сроком эксплуатации программного обеспечения такого класса: типовой жизненный цикл системы поддержки приёмной кампании ограничен изменениями нормативной базы приёма (Правила приёма в учреждения высшего образования пересматриваются раз в три-пять лет [\ref{src:laws:admission}]) и связанной с ними доработкой алгоритмов отбора.

Чувствительность результата к исходным параметрам расчёта невысока: снижение годовой экономии трудозатрат на 10~\% уменьшает чистый дисконтированный доход примерно на пять тысяч рублей, но не делает его отрицательным, а индекс рентабельности инвестиций остаётся выше единицы.

Таким образом, разработка системы компьютеризации ввода и мониторинга информации об абитуриентах при поступлении в БНТУ является экономически целесообразной: за четырёхлетний горизонт планирования инвестиции в её разработку окупаются с положительным чистым дисконтированным доходом, а рентабельность инвестиций превышает единицу.
Полученный экономический эффект достигается за счёт автоматизации всех ключевых этапов приёмной кампании БНТУ~--~от ввода данных абитуриента и его заявок до публикации результатов конкурсного отбора.
